\chapter{Linear Systems And Filters}
\label{ch:linear-systems-and-filters}

\section{Why Linear Systems Matter}

Chapter~\ref{ch:convolutions-for-sequences} ended with the key bridge:
a convolution kernel can be learned directly, or it can be generated by a recurrence.
This chapter explains that bridge more carefully.

The models in this chapter are linear. That may sound too weak for deep learning, but it
is exactly the right simplification. Deep sequence architectures often alternate between
linear temporal mixing, nonlinear pointwise transformations, normalization, and residual
connections. Attention is mostly a learned mixing operation inside a nonlinear block.
Convolution is a linear temporal filter inside a nonlinear block. S4 and S5 are also
linear temporal filters inside nonlinear blocks.

So the question is not whether a whole neural network should be linear. It should not.
The question is how much temporal structure can be captured by a powerful linear module
when that module is embedded in a deep nonlinear architecture.

\section{Scalar Linear Recurrences}

Start with the simplest possible recurrence:
\begin{equation}
  x_t = a x_{t-1} + b u_t,
  \qquad
  y_t = c x_t + d u_t,
  \label{eq:scalar-linear-recurrence}
\end{equation}
where $u_t$ is the input, $x_t$ is the state, and $y_t$ is the output. The parameters are
scalars $a,b,c,d$. Assume $x_0=0$ for now.

Unrolling gives
\[
\begin{aligned}
  x_1 &= b u_1, \\
  x_2 &= a b u_1 + b u_2, \\
  x_3 &= a^2 b u_1 + a b u_2 + b u_3.
\end{aligned}
\]
In general,
\[
  x_t = \sum_{i=0}^{t-1} a^i b u_{t-i}.
\]
Therefore
\[
  y_t
  =
  d u_t
  +
  c \sum_{i=0}^{t-1} a^i b u_{t-i}.
\]
This is already a convolution. Its kernel is
\[
  k_0 = cb + d,
  \qquad
  k_i = c a^i b \quad \text{for } i \geq 1,
\]
up to the convention for whether the direct term $d u_t$ is folded into $k_0$.

This scalar recurrence produces a geometric impulse response:
\[
  b,\; ab,\; a^2b,\; a^3b,\ldots.
\]
If $|a|<1$, the memory decays. If $|a|>1$, the response grows. If $a$ is close to $1$,
the memory decays slowly. This is the first appearance of a theme that will keep
returning: long memory requires dynamics that do not vanish too quickly, but stable
training requires dynamics that do not explode.

\section{Impulse Response}

The impulse response is the output produced by a unit input at one time and zero input
elsewhere. For the scalar recurrence above, let
\[
  u_1 = 1,
  \qquad
  u_t = 0 \text{ for } t > 1,
  \qquad
  x_0 = 0.
\]
Then
\[
  x_1=b,\quad x_2=ab,\quad x_3=a^2b,\quad \ldots
\]
and, ignoring the direct term for $t>1$,
\[
  y_t = c a^{t-1} b.
\]

For a finite CNN kernel, the impulse response becomes exactly zero after $K$ steps. For
the recurrence, the impulse response can continue forever:
\[
  c b,\; c a b,\; c a^2 b,\; c a^3 b,\ldots.
\]
This is an infinite impulse-response filter, usually abbreviated IIR. A finite
convolutional kernel is an FIR filter. Linear recurrences naturally give IIR filters.

In practice, a stable IIR response may become numerically negligible after some time, but
it is not forced to have a hard cutoff. This is one reason recurrences are attractive for
long-context modeling.

\section{Vector Linear Recurrences}

Now move from a scalar state to a vector state:
\begin{equation}
  \vx_t = \mA \vx_{t-1} + \mB \vu_t,
  \qquad
  \vy_t = \mC \vx_t + \mD \vu_t.
  \label{eq:discrete-linear-system}
\end{equation}
Here
\[
  \vx_t \in \R^N,\quad
  \vu_t \in \R^{C_{\mathrm{in}}},\quad
  \vy_t \in \R^{C_{\mathrm{out}}},
\]
and the parameter shapes are
\[
  \mA \in \R^{N \times N},\quad
  \mB \in \R^{N \times C_{\mathrm{in}}},\quad
  \mC \in \R^{C_{\mathrm{out}} \times N},\quad
  \mD \in \R^{C_{\mathrm{out}} \times C_{\mathrm{in}}}.
\]

The state dimension $N$ is not necessarily the same as the input channel count or the
model width. It is the size of the dynamical memory. In SSM papers this is often called
the state size.

Assuming $\vx_0=0$, unrolling gives
\[
\begin{aligned}
  \vx_t
  &=
  \mB \vu_t
  +
  \mA \mB \vu_{t-1}
  +
  \mA^2 \mB \vu_{t-2}
  + \cdots
  +
  \mA^{t-1}\mB \vu_1.
\end{aligned}
\]
Substituting into the output equation,
\begin{equation}
  \vy_t
  =
  \mD\vu_t
  +
  \sum_{i=0}^{t-1} \mC \mA^i \mB \vu_{t-i}.
  \label{eq:ssm-convolution-unrolled}
\end{equation}
So the recurrence defines a convolution kernel
\begin{equation}
  \mK_i = \mC \mA^i \mB.
  \label{eq:linear-system-kernel}
\end{equation}
This is the central formula for linear state space sequence models.

\section{Time Invariance}

The recurrence in \eqref{eq:discrete-linear-system} uses the same matrices
$\mA,\mB,\mC,\mD$ at every time step. Such a system is time invariant. If the same input
pattern occurs later in the sequence, the system responds the same way, shifted in time,
assuming the same initial state context.

Time invariance is what makes the convolution representation possible. The kernel
$\mK_i$ depends only on the lag $i$, not on the absolute time $t$. If the matrices changed
with time,
\[
  \vx_t = \mA_t \vx_{t-1} + \mB_t \vu_t,
\]
then the effect of $\vu_{t-i}$ on $\vy_t$ would depend on all the intermediate matrices.
There would no longer be a single fixed convolution kernel.

This is one major distinction between S4/S5 and Mamba. S4 and S5 are built around linear
time-invariant SSMs inside a neural block. Mamba introduces input-dependent parameters,
which breaks the fixed convolution view but adds content-dependent selectivity.

\section{Stability In Discrete Time}

To understand stability, ignore the input and direct output:
\[
  \vx_t = \mA \vx_{t-1}.
\]
Then
\[
  \vx_t = \mA^t \vx_0.
\]
The long-term behavior is controlled by powers of $\mA$. If these powers grow, the system
is unstable. If they decay, the system forgets its initial state.

For a scalar recurrence $x_t = a x_{t-1}$:
\[
  x_t = a^t x_0.
\]
The behavior is simple:
\[
\begin{array}{ll}
|a| < 1 & \text{decay}, \\
|a| = 1 & \text{persistent magnitude, possibly oscillatory}, \\
|a| > 1 & \text{growth}.
\end{array}
\]

For a matrix $\mA$, the analogous quantities are eigenvalues. If $\lambda$ is an
eigenvalue with eigenvector $\vv$, then
\[
  \mA \vv = \lambda \vv.
\]
Along that eigenvector direction,
\[
  \mA^t \vv = \lambda^t \vv.
\]
Thus, the magnitude $|\lambda|$ controls growth or decay of that mode. A discrete-time
linear system is stable, roughly, when all eigenvalues of $\mA$ lie inside the unit circle:
\[
  |\lambda_j(\mA)| < 1
  \quad \text{for all } j.
\]

This is the discrete-time version. Continuous-time systems use a different stability
condition: eigenvalues should have negative real part. Chapter~\ref{ch:state-space-models}
will connect these two views through discretization.

\section{Diagonalization And Modes}

Suppose $\mA$ can be diagonalized:
\[
  \mA = \mV \bm{\Lambda} \mV^{-1},
\]
where
\[
  \bm{\Lambda} = \diag(\lambda_1,\ldots,\lambda_N).
\]
Then
\[
  \mA^i = \mV \bm{\Lambda}^i \mV^{-1},
\]
and
\[
  \bm{\Lambda}^i = \diag(\lambda_1^i,\ldots,\lambda_N^i).
\]
This means the dynamics decompose into independent modes in the eigenbasis. Each mode
has its own decay rate and possible oscillation.

Using \eqref{eq:linear-system-kernel},
\[
  \mK_i
  =
  \mC \mV \bm{\Lambda}^i \mV^{-1} \mB.
\]
So the kernel is a mixture of terms like $\lambda_j^i$. In the scalar case, there was
only one geometric sequence. In the vector case, the impulse response is a learned mixture
of many geometric sequences.

This is a useful mental model for SSMs:
\[
  \text{SSM kernel} \approx \text{mixture of decays and oscillations}.
\]
Some modes can remember slowly. Some can respond quickly. Some can oscillate. The model
learns how input channels excite these modes through $\mB$ and how modes are read out
through $\mC$.

\section{Complex Eigenvalues And Oscillation}

Real matrices can have complex eigenvalues. For real-valued systems, complex eigenvalues
come in conjugate pairs:
\[
  \lambda = r e^{i\omega},
  \qquad
  \bar{\lambda} = r e^{-i\omega}.
\]
Powers of $\lambda$ behave as
\[
  \lambda^t = r^t e^{i\omega t}.
\]
The magnitude $r^t$ controls decay or growth. The angle $\omega t$ controls oscillation.

This is why complex-valued diagonal SSM parameterizations are natural. A mode with
$r$ close to $1$ and nonzero $\omega$ gives a slowly decaying oscillatory filter. Such
filters are useful for temporal data because real signals often have rhythms, periodic
components, and repeated temporal patterns.

The model output can still be real-valued. Complex conjugate modes combine to produce real
sines and cosines. In implementation, one may store complex parameters directly or use
equivalent real block matrices. The papers often use complex notation because it makes
the modal structure much cleaner.

\section{Vanishing, Exploding, And Long Memory}

The same eigenvalue logic explains the old RNN problem of vanishing and exploding
gradients. In a recurrent computation, forward signals and backward gradients both involve
repeated multiplication by transition-like matrices. Repeated powers can shrink or grow
rapidly.

If all eigenvalues have magnitudes much smaller than $1$, the system forgets quickly.
That is stable but short-memory. If eigenvalues have magnitudes larger than $1$, the
system can explode. If eigenvalues are very close to the unit circle, the system can keep
long memory, but optimization becomes delicate.

The goal is not simply to put every eigenvalue at $1$. That would preserve information
but could make the system brittle and poorly conditioned. The goal is to have a useful
spectrum of timescales:
\[
  \text{some fast modes} + \text{some medium modes} + \text{some slow modes}.
\]
HiPPO and S4-style initializations can be understood as principled attempts to initialize
such a spectrum for long-range sequence modeling.

\section{Transfer-Function Intuition}

Signal processing often describes linear systems by transfer functions. We do not need
the full machinery yet, but the intuition is helpful.

For a scalar convolution kernel $k_0,k_1,k_2,\ldots$, define the formal series
\[
  H(z) = \sum_{i=0}^{\infty} k_i z^{-i}.
\]
This says how the system weights different temporal frequencies. For a state space model,
the transfer function has the form
\[
  H(z)
  =
  \mC (z\mI - \mA)^{-1} \mB + \mD,
\]
up to convention choices about indexing. This compact expression contains the same
information as the impulse response $\mK_i = \mC\mA^i\mB$.

You do not need to manipulate transfer functions fluently to understand S4 or S5. The
important idea is that linear systems can be described in multiple equivalent ways:
\[
  \text{recurrence}
  \quad \leftrightarrow \quad
  \text{impulse response}
  \quad \leftrightarrow \quad
  \text{convolution}
  \quad \leftrightarrow \quad
  \text{frequency response}.
\]
Deep SSM papers move between these views because each one makes a different computational
or modeling property visible.

\section{Finite Versus Infinite Filters}

CNNs usually learn finite kernels:
\[
  \mK_0,\ldots,\mK_{K-1}.
\]
The parameter count grows linearly with kernel length $K$. If one wants a kernel that
can cover thousands of positions, direct learning becomes unattractive.

Linear systems generate kernels through matrix powers:
\[
  \mK_i = \mC\mA^i\mB.
\]
The parameter count depends on the state size and channel dimensions, not directly on the
sequence length. A compact state can generate a very long impulse response. This is the
basic compression advantage of a state space model.

There is also a constraint. Not every arbitrary long kernel can be represented well by a
small state. The SSM kernel must have structure: it is built from modes of $\mA$. That
structure is the inductive bias. It is helpful when the sequence really does contain
smooth, persistent, or oscillatory temporal dynamics. It may be limiting when the task
requires arbitrary content-dependent interactions.

\section{Multiple Inputs And Outputs}

In scalar examples, there is one input and one output. Real neural networks use many
channels. Equation \eqref{eq:discrete-linear-system} handles this with $\mB,\mC,\mD$.

The matrix $\mB$ decides how input channels drive the state:
\[
  \mB \vu_t \in \R^N.
\]
The matrix $\mC$ decides how the state is read into output channels:
\[
  \mC \vx_t \in \R^{C_{\mathrm{out}}}.
\]
The matrix $\mD$ gives an immediate input-to-output path:
\[
  \mD\vu_t.
\]

This distinction becomes important in S5. S4 often uses many single-input single-output
SSMs plus channel mixing around them. S5 emphasizes multi-input multi-output SSMs, where
one state space layer can jointly map multiple input channels to multiple output channels.

\section{What Is Linear, And What Is Learned?}

The recurrence in \eqref{eq:discrete-linear-system} is linear in the input sequence once
the parameters are fixed. But the parameters $\mA,\mB,\mC,\mD$ are learned by gradient
descent as part of a larger neural network.

In a deep block, the SSM may be surrounded by:
\[
  \text{normalization},\quad
  \text{linear projections},\quad
  \text{gates},\quad
  \text{nonlinearities},\quad
  \text{residual connections}.
\]
Thus the full model is nonlinear even if the temporal core is linear. This is exactly
analogous to a CNN: a convolution layer is linear, but a CNN block with activations and
residual paths is not.

\section{A Useful Mental Picture}

Think of the state $\vx_t$ as a bank of leaky memory cells. Each mode remembers the past
at a different rate, and possibly with a different oscillation. The input matrix $\mB$
decides which memory cells are written to. The transition matrix $\mA$ decides how those
memories decay, rotate, and mix. The output matrix $\mC$ decides which memories matter for
the current output.

This picture is not mathematically complete, but it is useful. S4 and S5 are mostly about
choosing a parameterization of this memory bank that is stable, expressive, and efficient
for long sequences.

\section{What To Remember}

\begin{itemize}
  \item A linear recurrence can be unrolled into a convolution.
  \item The impulse-response kernel of a discrete linear system is
  $\mK_i = \mC\mA^i\mB$.
  \item Eigenvalues of $\mA$ control decay, growth, and oscillation.
  \item Long memory requires modes that decay slowly but remain stable.
  \item Diagonalization turns matrix powers into independent modal powers
  $\lambda_j^i$.
  \item Complex eigenvalues naturally represent decaying oscillations.
  \item SSMs trade arbitrary learned kernels for compact, structured, dynamically
  generated kernels.
\end{itemize}

\section{Exercises}

\begin{enumerate}
  \item For the scalar recurrence $x_t = 0.9x_{t-1}+u_t$, $y_t=x_t$, write the first
  five impulse-response coefficients.
  \item Repeat the previous exercise with $a=-0.9$. How does the sign change affect the
  impulse response?
  \item Suppose $\mA=\diag(0.5,0.9,0.99)$. Which mode remembers longest? Which forgets
  fastest?
  \item Show by unrolling that the vector recurrence
  $\vx_t=\mA\vx_{t-1}+\mB\vu_t$ has terms
  $\mA^i\mB\vu_{t-i}$.
  \item Explain why making $\mA$ input-dependent would break the existence of a single
  fixed convolution kernel.
\end{enumerate}

\section{Prepares You For}

The next chapter formalizes state space models and explains how continuous-time dynamics
are discretized for neural network sequence layers. The new ingredient will be the
continuous system
\[
  \frac{d}{dt}\vx(t) = \mA\vx(t)+\mB\vu(t),
\]
and the question of how to convert it into a discrete recurrence suitable for sampled
sequence data.
